% GENERATED FILE. Edit canonical lessons, then run npm run generate:books.
% canonical-source-hash: fnv1a64:d3f5d7c39f065a16
% canonical-lessons: JA-C246-okoru, JA-C246-kenka, JA-C246-ayamaru, JA-C246-yurusu, JA-C246-zannen

\chapter{Anger and Apologies}
\label{ch:ja-anger-and-apologies}
\hlchaptermodality{246}

\begin{chapteropening}
\textbf{By the end of this chapter:} \emph{I can ask about quarrels and say who does not apologize.}

Complete the last lesson of chapter 246: I can ask about quarrels and say who does not apologize.
\end{chapteropening}

\section[\texorpdfstring{\ja{おこる} (okoru)}{おこる (okoru)}]{\ja{おこる} (okoru) — to get angry; to tell off}

\label{lesson:JA-C246-okoru}



\begin{quote}
Before the new one: say the Japanese for to laugh, then the Japanese for a smile.
\end{quote}

\subsection*{You'll want to know: \ja{おこる}}
\textbf{\ja{おこる}} — \emph{okoru} — \textquotedblleft{}to get angry; to tell off\textquotedblright{}.

\begin{quote}
\textbf{\ja{せんせいは} \ja{よく} \ja{おこりますか。}} — \emph{sensei wa yoku okorimasu ka} — \textquotedblleft{}Does the teacher often get angry?\textquotedblright{}
\end{quote}

Polite \textbf{\ja{おこります}}, \emph{okorimasu}; a godan verb.

\begin{cousinweb}
\textbf{\ja{おこる}} (\emph{okoru}) is a native Japanese word (wago), inherited rather than borrowed from Chinese; its kanji was attached later as a kun'yomi (JMdict, KANJIDIC2).
\end{cousinweb}

\begin{grammarlens}[title={What you've built}]
One more word for everyday life at A2.
\end{grammarlens}

\subsection*{Your turn}
\begin{itemize}
\raggedright
  \item \emph{Say it:} \emph{okoru}
  \item \emph{Say it:} \emph{okoru}, once more
  \item \emph{Say it:} \emph{sensei wa yoku okorimasu ka} — Does the teacher often get angry?
  \item \emph{Recall:} say the Japanese for to laugh, then the Japanese for a smile, then say \emph{okoru} again
\end{itemize}

\begin{tcolorbox}[breakable,colback=teal!4,colframe=teal!35!black,title={Before you move on}]
What does \ja{おこる} mean? (\textquotedblleft{}To get angry; to tell off\textquotedblright{}.) Say it once more.
\end{tcolorbox}
\section[\texorpdfstring{\ja{けんか} (kenka)}{けんか (kenka)}]{\ja{けんか} (kenka) — a quarrel; a fight}

\label{lesson:JA-C246-kenka}



\begin{quote}
Before the new one: say the Japanese for a smile, then the Japanese for to get angry.
\end{quote}

\subsection*{You'll want to know: \ja{けんか}}
\textbf{\ja{けんか}} — \emph{kenka} — \textquotedblleft{}a quarrel; a fight\textquotedblright{}.

\begin{quote}
\textbf{\ja{あには} \ja{ともだちと} \ja{けんかを} \ja{しません。}} — \emph{ani wa tomodachi to kenka o shimasen} — \textquotedblleft{}My elder brother does not quarrel with his friends.\textquotedblright{}
\end{quote}

\begin{cousinweb}
\textbf{\ja{けんか}} (\emph{kenka}) is Sino-Japanese (kango): two Chinese characters read with their on'yomi, readings borrowed from Chinese. They are \emph{ken}, \textquotedblleft{}noisy, boisterous\textquotedblright{} (Middle Chinese \emph{xjwon}), and \emph{ka}, \textquotedblleft{}noisy\textquotedblright{}, whose Middle Chinese \emph{xwae} is listed under a variant spelling that JMdict also gives (JMdict, KANJIDIC2; Middle Chinese in Baxter's transcription).
\end{cousinweb}

\begin{grammarlens}[title={What you've built}]
One more word for everyday life at A2.
\end{grammarlens}

\subsection*{Your turn}
\begin{itemize}
\raggedright
  \item \emph{Say it:} \emph{kenka}
  \item \emph{Say it:} \emph{kenka}, once more
  \item \emph{Say it:} \emph{ani wa tomodachi to kenka o shimasen} — My elder brother does not quarrel with his friends.
  \item \emph{Recall:} say the Japanese for a smile, then the Japanese for to get angry, then say \emph{kenka} again
\end{itemize}

\begin{tcolorbox}[breakable,colback=teal!4,colframe=teal!35!black,title={Before you move on}]
What does \ja{けんか} mean? (\textquotedblleft{}A quarrel; a fight\textquotedblright{}.) Say it once more.
\end{tcolorbox}
\section[\texorpdfstring{\ja{あやまる} (ayamaru)}{あやまる (ayamaru)}]{\ja{あやまる} (ayamaru) — to apologize}

\label{lesson:JA-C246-ayamaru}



\begin{quote}
Before the new one: say the Japanese for to get angry, then the Japanese for a quarrel.
\end{quote}

\subsection*{You'll want to know: \ja{あやまる}}
\textbf{\ja{あやまる}} — \emph{ayamaru} — \textquotedblleft{}to apologize\textquotedblright{}.

\begin{quote}
\textbf{\ja{おとうとは} \ja{いもうとに} \ja{あやまりますか。}} — \emph{otōto wa imōto ni ayamarimasu ka} — \textquotedblleft{}Does my younger brother apologize to my younger sister?\textquotedblright{}
\end{quote}

Polite \textbf{\ja{あやまります}}, \emph{ayamarimasu}; a godan verb.

\begin{cousinweb}
\textbf{\ja{あやまる}} (\emph{ayamaru}) is a native Japanese word (wago), inherited rather than borrowed from Chinese; its kanji was attached later as a kun'yomi (JMdict, KANJIDIC2).
\end{cousinweb}

\begin{grammarlens}[title={What you've built}]
One more word for everyday life at A2.
\end{grammarlens}

\subsection*{Your turn}
\begin{itemize}
\raggedright
  \item \emph{Say it:} \emph{ayamaru}
  \item \emph{Say it:} \emph{ayamaru}, once more
  \item \emph{Say it:} \emph{otōto wa imōto ni ayamarimasu ka} — Does my younger brother apologize to my younger sister?
  \item \emph{Recall:} say the Japanese for to get angry, then the Japanese for a quarrel, then say \emph{ayamaru} again
\end{itemize}

\begin{tcolorbox}[breakable,colback=teal!4,colframe=teal!35!black,title={Before you move on}]
What does \ja{あやまる} mean? (\textquotedblleft{}To apologize\textquotedblright{}.) Say it once more.
\end{tcolorbox}
\section[\texorpdfstring{\ja{ゆるす} (yurusu)}{ゆるす (yurusu)}]{\ja{ゆるす} (yurusu) — to forgive; to allow}

\label{lesson:JA-C246-yurusu}



\begin{quote}
Before the new one: say the Japanese for a quarrel, then the Japanese for to apologize.
\end{quote}

\subsection*{You'll want to know: \ja{ゆるす}}
\textbf{\ja{ゆるす}} — \emph{yurusu} — \textquotedblleft{}to forgive; to allow\textquotedblright{}.

\begin{quote}
\textbf{\ja{いもうとは} \ja{おとうとを} \ja{ゆるしますか。}} — \emph{imōto wa otōto o yurushimasu ka} — \textquotedblleft{}Does my younger sister forgive my younger brother?\textquotedblright{}
\end{quote}

Polite \textbf{\ja{ゆるします}}, \emph{yurushimasu}; a godan verb.

\begin{cousinweb}
\textbf{\ja{ゆるす}} (\emph{yurusu}) is a native Japanese word (wago), inherited rather than borrowed from Chinese; its kanji was attached later as a kun'yomi (JMdict, KANJIDIC2).
\end{cousinweb}

\begin{grammarlens}[title={What you've built}]
One more word for everyday life at A2.
\end{grammarlens}

\subsection*{Your turn}
\begin{itemize}
\raggedright
  \item \emph{Say it:} \emph{yurusu}
  \item \emph{Say it:} \emph{yurusu}, once more
  \item \emph{Say it:} \emph{imōto wa otōto o yurushimasu ka} — Does my younger sister forgive my younger brother?
  \item \emph{Recall:} say the Japanese for a quarrel, then the Japanese for to apologize, then say \emph{yurusu} again
\end{itemize}

\begin{tcolorbox}[breakable,colback=teal!4,colframe=teal!35!black,title={Before you move on}]
What does \ja{ゆるす} mean? (\textquotedblleft{}To forgive; to allow\textquotedblright{}.) Say it once more.
\end{tcolorbox}
\section[\texorpdfstring{\ja{ざんねん} (zannen)}{ざんねん (zannen)}]{\ja{ざんねん} (zannen) — a pity; regrettable}

\label{lesson:JA-C246-zannen}



\begin{quote}
Before the new one: say the Japanese for to apologize, then the Japanese for to forgive.
\end{quote}

\subsection*{You'll want to know: \ja{ざんねん}}
\textbf{\ja{ざんねん}} — \emph{zannen} — \textquotedblleft{}a pity; regrettable\textquotedblright{}.

\begin{quote}
\textbf{\ja{きのうの} \ja{しあいは} \ja{ざんねんでしたね。}} — \emph{kinō no shiai wa zannen deshita ne} — \textquotedblleft{}Yesterday's match was a pity, wasn't it.\textquotedblright{}
\end{quote}

\begin{cousinweb}
\textbf{\ja{ざんねん}} (\emph{zannen}) is Sino-Japanese (kango): two Chinese characters read with their on'yomi, readings borrowed from Chinese. They are \emph{zan}, \textquotedblleft{}remain\textquotedblright{} (Middle Chinese \emph{dzan}) and \emph{nen}, \textquotedblleft{}thought, feeling\textquotedblright{} (\emph{nemH}) (KANJIDIC2; Middle Chinese in Baxter's transcription).
\end{cousinweb}

\begin{grammarlens}[title={What you've built}]
One more word for everyday life at A2.
\end{grammarlens}

\subsection*{Your turn}
\begin{itemize}
\raggedright
  \item \emph{Say it:} \emph{zannen}
  \item \emph{Say it:} \emph{zannen}, once more
  \item \emph{Say it:} \emph{kinō no shiai wa zannen deshita ne} — Yesterday's match was a pity, wasn't it.
  \item \emph{Recall:} say the Japanese for to apologize, then the Japanese for to forgive, then say \emph{zannen} again
\end{itemize}

\begin{tcolorbox}[breakable,colback=teal!4,colframe=teal!35!black,title={Before you move on}]
What does \ja{ざんねん} mean? (\textquotedblleft{}A pity; regrettable\textquotedblright{}.) Say it once more.
\end{tcolorbox}
